\begin{figure}[H]
\centering
\resizebox{!}{0.9\textheight}{%
\begin{tikzpicture}[node distance=8mm and 20mm]

\node[term] (START) {Bắt đầu};
\node[proc, below=of START] (P1) {Nhận userId, menuPlanId};
\node[proc, below=of P1, text width=8.2cm] (P2) {Nạp MenuPlan kèm Items, Recipe,\\
RecipeIngredients với điều kiện\\
Id = menuPlanId VÀ UserId = userId};
\node[dec, below=of P2] (D1) {plan là null?};

\node[procs, right=32mm of D1, text width=3.4cm] (PERR) {Ném ngoại lệ: chặn truy cập chéo người dùng};
\node[term, below=of PERR] (ENDERR) {Kết thúc};

\node[proc, below=13mm of D1] (P3) {tatCa = danh sách rỗng};
\node[dec, below=of P3] (D2) {Còn suất ăn chưa duyệt?};

\node[procs, right=32mm of D2, text width=3.6cm] (P4) {Thêm toàn bộ RecipeIngredient của món ở suất này vào tatCa};

\node[proc, below=13mm of D2, text width=8.0cm] (P5) {Gộp nhóm tatCa theo khóa ghép (IngredientId, Unit)};
\node[proc, below=of P5, text width=8.0cm] (P6) {Với mỗi nhóm: tạo ShoppingListItem với tổng Quantity};
\node[dec, below=of P6] (D3) {Đã có ShoppingList của thực đơn này?};

\node[procs, right=32mm of D3, text width=3.4cm] (P7) {Xóa toàn bộ dòng cũ của danh sách};

\node[proc, below=13mm of D3] (P8) {Tạo ShoppingList mới};
\node[proc, below=of P8, text width=7.0cm] (P9) {Gán các dòng đã gộp và lưu};
\node[proc, below=of P9] (P10) {Trả về ShoppingList};
\node[term, below=of P10] (END) {Kết thúc};

\draw[fl] (START) -- (P1);
\draw[fl] (P1) -- (P2);
\draw[fl] (P2) -- (D1);
\draw[fl] (D1) -- (PERR) node[lbl, midway] {Có};
\draw[fl] (PERR) -- (ENDERR);
\draw[fl] (D1) -- (P3) node[lbl, midway] {Không};
\draw[fl] (P3) -- (D2);
\draw[fl] (D2) -- (P4) node[lbl, midway] {Có};
\draw[fl] (P4.south) -- ++(0,-4mm) -| (D2.west);
\draw[fl] (D2) -- (P5) node[lbl, midway] {Không};
\draw[fl] (P5) -- (P6);
\draw[fl] (P6) -- (D3);
\draw[fl] (D3) -- (P7) node[lbl, midway] {Có};
\draw[fl] (P7.south) -- ++(0,-6mm) -| (P9.east);
\draw[fl] (D3) -- (P8) node[lbl, midway] {Không};
\draw[fl] (P8) -- (P9);
\draw[fl] (P9) -- (P10);
\draw[fl] (P10) -- (END);

\end{tikzpicture}%
}
\caption{Lưu đồ thuật toán sinh danh sách đi chợ}
\label{fig:flow-di-cho}
\end{figure}
